% Class-compatibility preamble retained for reader excerpts of
%   Hiraku Nakajima, "Quiver Varieties and Branching", SIGMA 5 (2009) 003,
%   DOI 10.3842/SIGMA.2009.003, CC BY 4.0.
%
% The source pins the publisher class sigma, which is not part of the delivery baseline, so this
% file supplies the amsart baseline and then copies the author's own declaration block mechanically
% and unchanged: the aenume environment, the Theorem/Corollary/Lemma/Proposition and
% definition-style newtheorem declarations, the equation numbering, the author's reference macros
% (\lemref and the like), and the author's mathematical macros from the end of the preamble
% (\Gr, \Rcal, \dbar, \linf, \Irr, \Res, \Lareg, \MR, \aff and the rest).
%
% Two things are deliberately left out.  First, the publisher class internals and the single macro
% \rf that they serve, together with \graysquare and the colour \halfgray that go with it:
% nothing in the delivered excerpts uses them, and \rf depends on the internal dimension
% \y@boxdim of the sigma class.  Second, the author's package mathdots, which is not installed in
% the delivery TeX installation and whose commands no delivered excerpt uses.  Third, the author's
% no-op silencer declaration for the MR review numbers, because the delivery baseline package set
% already provides that command name and no delivered excerpt calls it.
%
% No mathematical macro, statement, symbol or assumption is introduced, changed or removed by this
% file.
\documentclass[10pt,a4paper]{amsart}
\usepackage{amssymb}
\usepackage{amsbsy,amscd,mathtools}
\usepackage{mathrsfs}
\usepackage{nicefrac}
\usepackage[all]{xy}
\usepackage{color}
\usepackage{hyperref}
\newenvironment{aenume}{%
  \begin{enumerate}%
  \renewcommand{\theenumi}{\alph{enumi}}%
  \renewcommand{\labelenumi}{\textup{(\theenumi)}}%
  }{\end{enumerate}}
%

\usepackage{verbatim}
\usepackage{version}
\usepackage{color}
\newenvironment{NB}{
\color{red}{\bf NB}. \footnotesize
}{}
% For a final version, NB disappears.
\excludeversion{NB}
% Older version
%\renewenvironment{NB}{%
%  \comment
%  }{\endcomment}
% The following is used inside NB
\newenvironment{NB2}{
\color{blue}{\bf NB}. \footnotesize
}{}
%       Theorem environments
%% \theoremstyle{plain} %% This is the default
\newtheorem{Theorem}[equation]{Theorem}
\newtheorem{Corollary}[equation]{Corollary}
\newtheorem{Lemma}[equation]{Lemma}
\newtheorem{Proposition}[equation]{Proposition}
%\newtheorem{Axiom}{Axiom}

\theoremstyle{definition}
\newtheorem{Definition}[equation]{Definition}
\newtheorem{Example}[equation]{Example}
\newtheorem{CounterExample}[equation]{Counter-Example}
\newtheorem{Notation}[equation]{Notation}
\newtheorem{Convention}[equation]{Convention}
\newtheorem{Conjecture}[equation]{Conjecture}

%\theoremstyle{remark}
\newtheorem{Remark}[equation]{Remark}
\newtheorem{Remarks}[equation]{Remarks}
%\newtheorem*{Claim}{Claim}
%\renewcommand{\theClaim}{} % to make the claim environment unnumbered
\newtheorem{Question}[equation]{Question}
%\renewcommand{\theQuestion}{}  % to make the notation environment
                               % unnumbered
\newtheorem{Exercise}[equation]{Exercise}
%\renewcommand{\theExercise}{}

\numberwithin{equation}{section}
%\numberwithin{Theorem}{section}
%\numberwithin{figure}{section}

\newcommand{\thmref}[1]{Theorem~\ref{#1}}
\newcommand{\secref}[1]{Section~\ref{#1}}
\newcommand{\lemref}[1]{Lemma~\ref{#1}}
\newcommand{\propref}[1]{Proposition~\ref{#1}}
\newcommand{\corref}[1]{Corollary~\ref{#1}}
\newcommand{\subsecref}[1]{Section~\ref{#1}}
\newcommand{\chapref}[1]{Chapter~\ref{#1}}
\newcommand{\defref}[1]{Definition~\ref{#1}}
\newcommand{\remref}[1]{Remark~\ref{#1}}
\newcommand{\remsref}[1]{Remarks~\ref{#1}}
%
%       Math definitions
%\newcommand{\lsp}[2]{{}^{#1}{#2}}
\newcommand{\lsp}[2]{{\mskip-.3mu}{}^{#1}\mskip-1mu{#2}}
\newcommand{\defeq}{\overset{\operatorname{\scriptstyle def}}{=}}
\newcommand{\interior}{\rule{6pt}{.4pt}\rule{.4pt}{6pt}} % interior
                                                         % product
\newcommand{\C}{{\mathbb C}}
\newcommand{\Z}{{\mathbb Z}}
\newcommand{\Q}{{\mathbb Q}}
\newcommand{\R}{{\mathbb R}}
% Projective spaces
\newcommand{\proj}{{\mathbb P}}
%\newcommand{\CP}{\operatorname{\C P}}
\newcommand{\CP}{\proj}
\newcommand{\RP}{\operatorname{\R P}}
% Lie Groups
\newcommand{\SL}{\operatorname{\rm SL}}
\newcommand{\SU}{\operatorname{\rm SU}}
\newcommand{\GL}{\operatorname{GL}}
\newcommand{\PGL}{\operatorname{PGL}}
\newcommand{\U}{\operatorname{\rm U}}
\newcommand{\SO}{\operatorname{\rm SO}}
\newcommand{\grpSp}{\operatorname{\rm Sp}}
% Lie algebras
\newcommand{\algsl}{\operatorname{\mathfrak{sl}}} % because \sl="slant"
\newcommand{\su}{\operatorname{\mathfrak{su}}}
\newcommand{\gl}{\operatorname{\mathfrak{gl}}}
\newcommand{\algu}{\operatorname{\mathfrak{u}}}
\newcommand{\so}{\operatorname{\mathfrak{so}}}
\newcommand{\algsp}{\operatorname{\mathfrak{sp}}}
\newcommand{\g}{{\mathfrak g}}
\newcommand{\h}{{\mathfrak h}}
% Useful symbols
\newcommand{\Spec}{\operatorname{Spec}\nolimits}
\newcommand{\Proj}{\operatorname{Proj}\nolimits}
\newcommand{\End}{\operatorname{End}}
\newcommand{\Hom}{\operatorname{Hom}}
\newcommand{\Ext}{\operatorname{Ext}}
\newcommand{\Ker}{\operatorname{Ker}}
\newcommand{\Coker}{\operatorname{Coker}}
\newcommand{\Ima}{\operatorname{Im}}
\newcommand{\coker}{\operatorname{coker}}
\newcommand{\Endskew}{\operatorname{Endskew}}
\newcommand{\im}{\mathop{\text{\rm im}}\nolimits}
%\newcommand{\dim}{\mathop{\text{\rm dim}}\nolimits}
\newcommand{\rank}{\operatorname{rank}}
\newcommand{\ind}{\mathop{\text{\rm ind}}\nolimits}
\newcommand{\tr}{\operatorname{tr}}
\newcommand{\Tr}{\operatorname{Tr}}
\newcommand{\DB}{\overline{\partial}}
\newcommand{\pd}[2]{\frac{\partial#1}{\partial#2}}
\newcommand{\deriv}[1]{\frac{d}{d#1}}
\newcommand{\id}{\operatorname{id}}
\newcommand{\ve}{\varepsilon}
%
% notations used in this paper
%
% \newcommand{\vin}[1]{\operatorname{in}(#1)} % incoming vertex
% \newcommand{\vout}[1]{\operatorname{out}(#1)} % outgoing vertex
\newcommand{\vin}[1]{\operatorname{i}(#1)} % incoming vertex
\newcommand{\vout}[1]{\operatorname{o}(#1)} % outgoing vertex
\newcommand{\bM}{{\mathbf M}} % the vector space
\newcommand{\bN}{{\mathbf N}} % the vector space
\newcommand{\M}{{\mathfrak M}} % the moduli space
\newcommand{\Mreg}{\M^{\operatorname{s}}}
\newcommand{\La}{{\mathfrak L}} % the Lagrangian variety
\newcommand{\Fi}{{\mathfrak F}} % the fixed point variety
\newcommand{\dslash}{/\!\!/} % for algebro-geometric quotient (double slash)
\newcommand{\bv}{{\mathbf v}} % the vector v
\newcommand{\bw}{{\mathbf w}} % the vector w
\newcommand{\bC}{{\mathbf C}} % the Cartan matrix
\newcommand{\bA}{{\mathbf A}} % the adjacency matrix
\newcommand{\bI}{{\mathbf I}} % the identity matrix
\newcommand{\bn}{{\mathbf n}} % the kernel of the extended Cartan matrix
\newcommand{\bu}{{\mathbf u}}
\newcommand{\be}{{\mathbf e}}
\newcommand{\bq}{{\mathbf q}}
\newcommand{\bs}{\mathbf s}
\newcommand{\bc}{\mathbf c}
\newcommand{\codim}{\operatorname{codim}} % codimension
\newcommand{\gr}{\operatorname{gr}} % grade
\newcommand{\topdeg}{{\operatorname{top}}} % top degree
\newcommand{\Uq}{{\mathbf U}_q(\mathfrak g)} % the QUE algebra
\newcommand{\Un}{{\mathbf U}_q}
\newcommand{\Um}{\Tilde{\mathbf U}_q(\mathfrak g)} % the modified
                                                 % enveloping algebra
\newcommand{\Ua}{{\mathbf U}_q(\widehat{\mathfrak g})} % the quantum affine
                                                 % algebra
\newcommand{\Uae}{{\mathbf U}_{\varepsilon}(\widehat{\mathfrak g})} %
                                                 % the quantum affine
                                                 % algebra
\newcommand{\Ul}{{\mathbf U}_q({\mathbf L}{\mathfrak g})} % the quantum
                                                      % loop algebra
\newcommand{\Uma}{\Tilde{\mathbf U}^\bw_q(\widehat{\mathfrak g})}
%                                                 the quantum affine algebra
\newcommand{\Ui}{{\mathbf U}^{\Z}_q(\mathfrak g)}
\newcommand{\Ue}{{\mathbf U}_{\varepsilon}(\mathfrak g)}
\newcommand{\Ueone}{{\mathbf U}_{\varepsilon = 1}(\mathfrak g)}
\newcommand{\Uli}{{\mathbf U}^{\Z}_q({\mathbf L}\mathfrak g)}
\newcommand{\Ule}{{\mathbf U}_{\varepsilon}({\mathbf L}\mathfrak g)}
\newcommand{\bU}{\mathbf U} % the enveloping algebra
\newcommand{\Uqm}{{\widetilde{\mathbf U}_q}} % its modified form
\newcommand{\Pa}{{\mathfrak P}} % Hecke correspondence
\newcommand{\HomE}{\operatorname{E}}
\newcommand{\HomL}{\operatorname{L}}
\newcommand{\Zreg}{Z^{\operatorname{reg}\nolimits}}
\newcommand{\hiw}{{L}} % the highest weight module
%      \interval is used to provide better spacing after a [ that
%      is used as a closing delimiter.
\newcommand{\interval}[1]{\mathinner{#1}}
\newcommand{\shfO}{\mathcal O}
\newcommand{\Mp}{\widetilde\M}
\newcommand{\Mpp}{\widehat\M}
\newcommand{\Wedge}{{\textstyle \bigwedge}}
%
\newcommand{\Zw}{Z(\bw)}
\newcommand{\Mw}{\M(\bw)}
\newcommand{\Law}{\La(\bw)}
%\newcommand{\Zw}{Z(\bullet,\bullet;\bw)}
%\newcommand{\Mw}{\M(\bullet,\bw)}
%\newcommand{\Law}{\La(\bullet,\bw)}
\newcommand{\ngpiso}{\overset{.}{=}}
\newcommand{\idl}{\mathfrak I}
\newcommand{\Ktop}{K_{\operatorname{top}}}
\newcommand{\Kztop}{K_{0,\operatorname{top}}}
\newcommand{\Kotop}{K_{1,\operatorname{top}}}

\newcommand{\qH}{\mathbb H}
\newcommand{\D}{\mathcal D}
\newcommand{\A}{\mathcal A}
\newcommand{\incarrow}{\xrightarrow{\mathrm{inclusion}}}
\newcommand{\proarrow}{\xrightarrow{\mathrm{projection}}}
%\newcommand{\oV}[1]{\overset{\circ}{V_{#1}}}
\newcommand{\oV}[1]{\breve{V}_{#1}}
\newcommand{\bV}[2]{\bar{V}^{#1}_{#2}}
\newcommand{\bB}{\bar{B}}
\newcommand{\bi}{\bar{a}}
\newcommand{\bj}{\bar{b}}
\newcommand{\bphi}{\bar{\phi}}
\newcommand{\Gr}{\operatorname{Gr}}
\newcommand{\Pp}{\mathbf P}
\newcommand{\Rcal}{\mathcal R}
\newcommand{\dbar}{\overline\partial}
\newcommand{\linf}{\ell_\infty}
\newcommand{\Irr}{\operatorname{Irr}}
\newcommand{\Res}{\operatorname{Res}}
\newcommand{\Lareg}{\mathfrak L^{\operatorname{s}}}
%\renewcommand{\MR}[1]{}
\newcommand{\aff}{\mathrm{af\/f}}
\newcommand{\affLevi}{(\g_{I_0^0})_{\aff}}


\newcommand{\hf}{\hfill}

\newcommand{\graysquare}{\color{halfgray}\blacksquare}

%\setcounter{tocdepth}{1}

%\usepackage[bookmarks=false]{hyperref}
%\usepackage[dvips,bookmarks=false]{hyperref}
%\usepackage[initials]{amsrefs}

\usepackage{cleveref}
\crefname{originallemma}{original lemma}{original lemmas}
\Crefname{originallemma}{Original lemma}{Original lemmas}
\makeatletter
\newlabel{lem:root}{{2.10}{source}{}{sourcepointer-1}{}}
\newlabel{lem:root@cref}{{[originallemma][0][]2.10}{[1][0][]source}{}{}{}}
\makeatother
\title{A criterion for the non-emptiness of the regular part of a quiver variety in terms of roots and weights}
\author{Hiraku Nakajima}
\date{doi:10.3842/SIGMA.2009.003; CC BY 4.0}
\begin{document}
\maketitle
\noindent\textit{Editorial note.} All delivered mathematical text is the author's own, copied from the pinned publisher source sources/algebra-fl01-quiver-varieties-and-branching/pinned-original-source.tex; no statement or proof step was rewritten, reordered or removed.\par\medskip
\noindent\textit{Editorial note.} Pairing provenance: the counted statement is the author's Theorem environment at source lines 1219-1243 (label thm:CB) and the counted proof is the author's proof environment at source lines 1253-1282. The two are not adjacent: the author's unlabelled Remark at source lines 1245-1251 lies between them and is delivered by neither span. The pairing is established by the proof's own content, whose two parts (1) and (2) prove exactly the two numbered parts of the statement, and by the remark's status as a conjecture about more general parameters rather than a statement with its own proof.\par\medskip
\noindent\textit{Editorial note.} Preamble provenance: the source pins the publisher class sigma, which the delivery wrapper does not accept, so this card uses sources/algebra-fl01-quiver-varieties-and-branching/retained-preamble.tex, which supplies the amsart baseline and copies the author's declaration block of pinned source lines 43-239 unchanged. No macro was invented and no mathematical symbol was changed.\par\medskip
\noindent\textit{Editorial note.} Bibliography provenance: the counted proof cites Theorem 1.2 of the author's reference CB. That citation is resolved with sources/algebra-fl01-quiver-varieties-and-branching/bibliography.tex, an excerpt of the author's own thebibliography environment (pinned source lines 4098-4257) copied character for character with all thirty-three entries kept; the entry CB is present verbatim. No entry was rewritten and no citation key was changed.\par\medskip
\noindent\textit{Editorial note.} External references are navigation pointers only; every deliberately kept step of the author's argument stays inside the counted spans.\par\medskip
\noindent\textit{Editorial note.} Selection provenance: this card was opened under the file-based control contract (CONTROL.md), with the mechanical pre-card gate run on its statement and proof interval before the card was written and reporting the interval free.\par\medskip
\noindent\textit{Allowed standard prerequisites (editorial).}\begin{itemize}
\item The moduli space M\_zeta(v,w) of stable modules, its regular part, and the stability parameter zeta with the Cartan form pairing p(x) = 1 - (1/2) x-transpose C x. Set up in the author's Sections 1 and 2, before the counted statement. Fixes every object of the two numbered parts, including the quantity p and the roots appearing in the inequalities.
\item The classification of simple modules via the cited theorem of CB, and the hyper-Kahler moment map description of stable modules. Used inside the counted proof at source lines 1254-1259, where the author invokes Theorem 1.2 of the cited paper CB. Supplies the criterion for the existence of a simple module with prescribed moment map, which is the whole content of part (1).
\item The construction of the extended quiver and the root-space identification for the extended graph. The author's lemma labelled lem:root (source lines 1184-1196), which the counted proof invokes. Provides the dimension-vector and root bookkeeping used in part (2).
\end{itemize}\medskip
\ifdefined\hypertarget\hypertarget{sourcepointer-1}{}\fi
\noindent\textit{Source reference.} Original lemma 2.10: sources/algebra-fl01-quiver-varieties-and-branching/pinned-original-source.tex, source lines 1184--1184. Navigation pointer only: the counted proof invokes the author's root-space lemma in its part (2); that lemma has its own card in this lane and its text is not reproduced here.\par\medskip
\begin{Theorem}\label{thm:CB}
  \textup{(1)} Suppose $W = 0$ and consider a dimension vector $\bv$
  with $\zeta\cdot \bv = 0$. Then $\Mreg_\zeta(\bv,0)\neq\varnothing$ if
  and only if the following holds\textup:
\begin{itemize}\itemsep=0pt
\item $\bv$ is a positive root, and $p(\bv) > \sum_{t=1}^r
  p(\beta^{(t)})$ for any decomposition $\bv = \sum_{t=1}^r
  \beta^{(t)}$ with $r \ge 2$ and $\beta^{(t)}$ a positive root with
  $\zeta\cdot \beta^{(t)} = 0$ for all $t$,
\end{itemize}
where $p(x) = 1 - \frac12 {}^t x \bC x$.

\textup{(2)} Suppose $\bw\neq 0$. Then
$\Mreg_\zeta(\bv,\bw)\neq\varnothing$ if and only if the following
holds:
\begin{itemize}\itemsep=0pt
\item $\bw - \bv$ is a weight of the integrable highest weight
  representation $V(\bw)$ of the highest weight $\bw$, and
  $\lsp{t}{\bv}(\bw - \frac12 \bC \bv) > \lsp{t}{\bv^0}(\bw - \frac12
  \bC \bv^0) + \sum_{t=1}^r p(\beta^{(t)})$ for any decomposition $\bv
  = \bv^0 + \sum_{t=1}^r \beta^{(t)}$ with $r \ge 1$, $\bw - \bv^0$ is
  a weight of $V(\bw)$, and $\beta^{(t)}$ a positive root with
  $\zeta\cdot \beta^{(t)} = 0$ for all $t$.
\end{itemize}
\end{Theorem}

\begin{proof}
(1) A $\zeta$-stable module with $\mu = 0$ corresponds to a point
satisfying the hyper-K\"ahler moment map $(\mu_\R,\mu) = (\zeta,0)$,
which correspond a $0$-stable module (i.e., a simple module) with $\mu
= \zeta$ after a rotation of the complex structures. Then by
\cite[Theorem~1.2]{CB} such a simple module exists if and only if the
above condition holds.

(2) For the given $\bw$ we construct the quiver $(\Tilde{I},\Tilde{E})$ as
above. We consider the dimension vector $\Tilde{\bv} = \bv + \alpha_\infty$
and def\/ine the stability parameter $\Tilde{\zeta} = (\zeta,\zeta_\infty)$ by
\(
\zeta\cdot \dim \bv + \zeta_\infty = 0.
\)
Then $\Mreg_\zeta(\bv,\bw)$ is isomorphic to the quiver variety
$\Mreg_{\Tilde{\zeta}}(\Tilde{\bv},0)$ associated with
$(\Tilde{I},\Tilde{E})$ \cite[the end of Introduction]{CB}.
%
Therefore we can apply the criterion in (1) together with the
observation \lemref{lem:root}. We f\/irst note that
\[
  p(\Tilde{\bv}) = 1 - \tfrac12 \lsp{t}{\Tilde{\bv}} \Tilde{\bC} \Tilde{\bv}
  = \lsp{t}{\bv}\left(\bw - \tfrac12 \bC \bv\right).
\]
Also if we have a decomposition $\Tilde{\bv} = \sum \beta^{(t)}$, one
of $\beta^{(t)}$ has $1$ in the entry $\infty$, and other
$\beta^{(t)}$'s have $0$ in the entry $\infty$. We rewrite the former
as $\bv^0 + \alpha_\infty$, and identify the latter with positive
roots for $(I,E)$. Now the assertion follows from (1).
\end{proof}

%% Bibliography excerpt retained for reader excerpts of
%   Hiraku Nakajima, "Quiver Varieties and Branching",
%   SIGMA 5 (2009) 003, DOI 10.3842/SIGMA.2009.003.
%%
%% Every entry below is copied character for character from the author's own
% thebibliography environment of the pinned publisher source (lines 4098-4257).
%% No entry was rewritten, reordered, merged, split or invented, and no citation
% key was changed.  All 33 entries are kept so that any citation of a delivered
%% excerpt resolves without guessing an entry.
\begin{thebibliography}{99}


\footnotesize\itemsep=0pt

\bibitem{Bando}
Bando S.,  Einstein--Hermitian metrics on noncompact  K \"ahler
  manifolds, in Einstein Metrics and Yang--Mills connections (Sanda, 1990),
 {\it  Lecture Notes in Pure and Appl. Math.}, Vol.~145, Dekker, New York, 1993,
27--33. %\MR{MR1215276 (94d:32040)}

\bibitem{Baranovsky}
Baranovsky V.,  Moduli of sheaves on surfaces and action of the oscillator
  algebra, {\it J. Differential Geom.} \textbf{55} (2000),  193--227, \href{http://arxiv.org/abs/math.AG/9811092}{math.AG/9811092}.
  %\MR{MR1847311 (2002h:14071)}

\bibitem{BM}
Borho W., MacPherson R.,  Partial resolutions of nilpotent varieties,
  in Analysis and Topology on Singular Spaces, II, III (Luminy, 1981),
 {\it  Ast\'erisque}, Vol.~101, Soc. Math. France, Paris, 1983,  23--74.
  %\MR{MR737927 (85j:14087)}

\bibitem{braverman-2007}
Braverman A., Finkelberg M.,  Pursuing the double af\/f\/ine {G}rassmannian
  {I}: transversal slices via instantons on ${A}_k$-singularities, \href{http://arxiv.org/abs/0711.2083}{arXiv:0711.2083}.

\bibitem{braverman-prep}
Braverman A., Finkelberg M., Private communication, 2008.

\bibitem{bk}
Braverman A., Kazhdan D., The spherical {H}ecke algebra for af\/f\/ine
  {K}ac--{M}oody groups {I}, \href{http://arxiv.org/abs/0809.1461}{arXiv:0809.1461}.

\bibitem{CG}
Chriss N., Ginzburg V.,  Representation theory and complex geometry,
  Birkh\"auser Boston Inc., Boston, MA, 1997. %\MR{MR1433132 (98i:22021)}

\bibitem{CB}
Crawley-Boevey W., Geometry of the moment map for representations of
  quivers, {\it Compositio Math.} \textbf{126} (2001),   257--293.
 % \MR{MR1834739 (2002g:16021)}

\bibitem{MR675108}
 Frenkel I.B., Representations of af\/f\/ine {L}ie algebras, {H}ecke modular
  forms and {K}orteweg--de~{V}ries type equations, in Lie Algebras and Related
  Topics ({N}ew {B}runswick, {N}.{J}., 1981), {\it Lecture Notes in Math.}, Vol.~933,
  Springer, Berlin~-- New York, 1982, 71--110. %\MR{MR675108 (84f:17006)}

\bibitem{Grojnowski}
Grojnowski I., Instantons and af\/f\/ine algebras. I.~The Hilbert
  scheme and vertex operators, {\it Math. Res. Lett.} \textbf{3} (1996),  275--291,
  \href{http://arxiv.org/abs/alg-geom/9506020}{alg-geom/9506020}. %\MR{MR1386846 (97f:14041)}

\bibitem{Hasegawa}
Hasegawa K., Spin module versions of Weyl's reciprocity theorem for
  classical Kac--Moody Lie algebras -- an application to branching rule
  duality, {\it Publ. Res. Inst. Math. Sci.} \textbf{25} (1989),  741--828.
 % \MR{MR1031225 (91f:17023)}

\bibitem{Kac}
 Kac V.G., Inf\/inite-dimensional Lie algebras, 3rd ed., Cambridge
  University Press, Cambridge, 1990. %\MR{MR1104219 (92k:17038)}

\bibitem{KS}
Kashiwara M., Saito Y., Geometric construction of crystal bases, {\it Duke
  Math. J.} \textbf{89} (1997),  9--36,  \mbox{\href{http://arxiv.org/abs/q-alg/9606009}{q-alg/9606009}}. %\MR{MR1458969 (99e:17025)}

\bibitem{King}
 King A.D.,  Moduli of representations of f\/inite-dimensional algebras,
 {\it Quart. J. Math. Oxford Ser. (2)} \textbf{45} (1994),  515--530.
 % \MR{MR1315461 (96a:16009)}

\bibitem{Kr}
 Kronheimer P.B.,  The construction of {ALE} spaces as hyper-{K}\"ahler
  quotients, {\it J. Differential Geom.} \textbf{29} (1989),    665--683.
 % \MR{MR992334 (90d:53055)}

\bibitem{KN}
 Kronheimer P.B., Nakajima H., Yang--Mills instantons on {ALE}
  gravitational instantons, {\it Math. Ann.} \textbf{288} (1990),   263--307.
  %\MR{MR1075769 (92e:58038)}

\bibitem{Lu-Green}
Lusztig G.,  Green polynomials and singularities of unipotent classes,
 {\it  Adv. in Math.} \textbf{42} (1981),  169--178. %\MR{MR641425 (83c:20059)}

\bibitem{Lu-crystal}
Lusztig G.,  Canonical bases arising from quantized enveloping algebras.
  II, in Common
  Trends in Mathematics and Quantum Field Theories (Kyoto, 1990),
  {\it Progr. Theoret. Phys. Suppl.} (1990), no.~102, 175--201. %\MR{MR1182165   (93g:17019)}

\bibitem{Malkin}
Malkin A.,  Tensor product varieties and crystals: the {$ADE$} case, {\it Duke
  Math. J.} \textbf{116} (2003),  477--524, \href{http://arxiv.org/abs/math.AG/0103025}{math.AG/0103025}. %\MR{MR1958096 (2004e:17003)}

\bibitem{MV}
Mirkovi{\'c} I., Vilonen K.,  Perverse sheaves on af\/f\/ine
  {G}rassmannians and {L}anglands duality, {\it Math. Res. Lett.} \textbf{7} (2000),
 13--24, \href{http://arxiv.org/abs/math.AG/9911050}{math.AG/9911050}. %\MR{MR1748284 (2001h:14020)}

\bibitem{Nagao-2007}
Nagao K.,  Quiver varieties and Frenkel--Kac construction,   \href{http://arxiv.org/abs/math.RT/0703107}{math.RT/0703107}.

\bibitem{Na-quiver}
Nakajima H., Instantons on {ALE} spaces, quiver varieties, and
  Kac--Moody algebras, {\it Duke Math. J.} \textbf{76} (1994), 365--416.
  %\MR{MR1302318 (95i:53051)}

\bibitem{Na-proc}
Nakajima H., Varieties associated with quivers, in Representation Theory of
  Algebras and Related Topics (Mexico City, 1994), {\it CMS Conf. Proc.}, Vol.~19,
  Amer. Math. Soc., Providence, RI, 1996, 139--157.
  %\MR{MR1388562   (97m:16022)}

\bibitem{Na-alg}
Nakajima H., Quiver varieties and Kac--Moody algebras, {\it Duke Math. J.}
  \textbf{91} (1998), 515--560. %\MR{MR1604167 (99b:17033)}

\bibitem{Lecture}
Nakajima H., Lectures on Hilbert schemes of points on surfaces, {\it University
  Lecture Series}, Vol.~18, American Mathematical Society, Providence, RI, 1999.
  %\MR{MR1711344 (2001b:14007)}

\bibitem{Na-qaff}
Nakajima H., Quiver varieties and f\/inite-dimensional representations of
  quantum af\/f\/ine algebras, {\it J. Amer. Math. Soc.} \textbf{14} (2001),
  145--238, \href{http://arxiv.org/abs/math.QA/9912158}{math.QA/9912158}. %\MR{MR1808477 (2002i:17023)}

\bibitem{Na-tensor}
Nakajima H., Quiver varieties and tensor products, {\it Invent. Math.}
  \textbf{146} (2001), 399--449, \href{http://arxiv.org/abs/math.QA/0103008}{math.QA/0103008}. %\MR{MR1865400 (2003e:17023)}

\bibitem{MR1989196}
Nakajima H., Geometric construction of representations of af\/f\/ine algebras,
  in Proceedings of the International Congress of Mathematicians, Vol. I (Beijing,
  2002), Higher Ed. Press, Beijing, 2002, 423--438, \href{http://arxiv.org/abs/math.QA/0212401}{math.QA/0212401}. %\MR{MR1989196 (2004e:17012)}

\bibitem{Na-reflect}
Nakajima H.,  Ref\/lection functors for quiver varieties and {W}eyl group
  actions, {\it Math. Ann.} \textbf{327} (2003), 671--721.
  %\MR{MR2023313   (2004k:16036)}

\bibitem{Na-ADHM}
Nakajima H.,  Sheaves on $ALE$ spaces and quiver varieties, {\it Mosc. Math. J.}
  \textbf{7} (2007),  699--722.

\bibitem{NT}
Nakanishi T., Tsuchiya A.,  Level-rank duality of {WZW} models in
  conformal f\/ield theory, {\it Comm. Math. Phys.} \textbf{144} (1992),
  351--372. %\MR{MR1152377 (93a:81181)}

\bibitem{Rudakov}
Rudakov A.,  Stability for an abelian category, {\it J. Algebra} \textbf{197}
  (1997),   231--245. %\MR{MR1480783 (98j:14011)}

\bibitem{Saito}
Saito Y., Crystal bases and quiver varieties, {\it Math. Ann.} \textbf{324}
  (2002),  675--688, \href{http://arxiv.org/abs/math.QA/0111232}{math.QA/0111232}. %\MR{MR1942245 (2004a:17023)}


\end{thebibliography}
\end{document}
